% Independent derivation. Prefix ZIL. Requires amsmath, amssymb, hyperref.
\section{Supported-zero inversion, the full lattice, and retained amplitudes}
\label{sec:zero-inversion-full-lattice-review}

\paragraph{Sources and the claim being tested.}
The user-supplied source \nolinkurl{globalization_note.tex},
\emph{A Note on the Globalization Semiring $G(\mathbb Z)$}, has an empty
author field and the date March 2026. Its entire source was read for
this calculation. The exact source has SHA--256
\begin{center}\small
\nolinkurl{ce06ac01c32e0452fe36121ac0f90e4eff8fad469b1cda7a6cf2d73116077a02}.
\end{center}
The relevant source locators are:
\begin{itemize}
\item Proposition \texttt{prop:units-irr}, lines 282--300,
and Corollary \texttt{cor:zeta}, lines 406--430;
\item Theorem \texttt{thm:localization-away-from-zero}, lines 436--480;
\item Theorem \texttt{thm:frac-collapse}, lines 482--534;
\item the definition at lines 817--825, followed by Theorem
\texttt{thm:class-semigroup} and its proof at lines 827--871.
\end{itemize}
The source's bibliography names Golan, Deitmar, Lorscheid and Neukirch;
those bibliography entries are not represented here as author sources
independently read during this calculation. The note's empty author
field does not authorize assigning it a human author.

The full-lattice carrier and its supported-zero localization are
existing programme results. Their public proof is
\href{https://github.com/KokunoYumeto/zeta-function-research-reader/blob/3c78cfbd9e194d35117b33277c421c6564ce3dbf/workbenches/splitzero-tandem/continuations/20260923-full-prime-support/031/01_REBUILT_MATHEMATICS.tex}{031,
Theorems \texttt{thm:rebuilt-full-support} and \texttt{thm:e-stalk},
equation \texttt{eq:e-element-localization}}.
The earlier foundation is retained under its byline The Clankers in
\href{https://github.com/KokunoYumeto/zeta-function-research-reader/blob/a2851f9eb352e7e4376bc629de86fa4ed9c1c434/workbenches/splitzero-tandem/foundations/split-support-geometry-v11/split_support_geometry_arithmetic_curve_v11.tex}{\emph{Split
Support Geometry: Universal Zero Fibres, Arithmetic Curves, Frobenius
and Perfect Quantization}, v11, Definition \texttt{def:lattice-split}}.
The user's attribution of the earlier scalar supported-zero prime
observation to Gemini Pro 2.5 is retained as user-reported provenance;
this calculation does not claim that observation as new.

The precise characteristic-one comparison being used is ASL1--22 of
the complete included arithmetic-site lattice-map proof, read in full
through ASL22. Its source is
\begin{center}\small
\nolinkurl{ARITHMETIC_SITE_LATTICE_MAP.tex},\\
SHA--256
\nolinkurl{a03b6c40ba06653fe3a481592170d1c4c43ff16d1b4af53959b3e0ab1fe0b492}.
\end{center}
ASL identifies its human framework source as Alain Connes and Caterina
Consani, \href{https://arxiv.org/abs/1502.05580}{\emph{Geometry of the
Arithmetic Site}}; the following proofs use the displayed ASL
constructions and give the required algebra independently. They do
not claim a new reading of that entire human paper.

The statement to be tested is the user-supplied wording
``zero is inmvertable this means rh is wrong''. We calculate exactly
which element is inverted, what its inverse is, which identities and
prime points result, and what information survives in the actual
joint coefficient map. No analytic zero is silently identified with
the semiring's absorbing zero.

\subsection{The original two zeros and a general idempotent localization}

Let $R$ be a nonzero commutative unital ring and $L$ a nontrivial
bounded distributive lattice. Its semiring operations are join and
meet. Retain the complete carrier
\begin{equation}\tag{ZIL1}
\begin{gathered}
 S=G_L(R)=\{(0_R,\lambda):\lambda\in L\}
                 \cup\{(r,1_L):r\in R\},\\
 (r,\lambda)+(s,\nu)=(r+s,\lambda\vee\nu),\qquad
 (r,\lambda)(s,\nu)=(rs,\lambda\wedge\nu),\\
 z_\lambda=(0_R,\lambda),\quad
 \tau=z_{0_L}=0_S,\quad e=z_{1_L},\quad
 1_S=(1_R,1_L),\\
 q:S\longrightarrow R,\quad q(r,\lambda)=r,\qquad
 \sigma:S\longrightarrow L,\quad\sigma(r,\lambda)=\lambda.
\end{gathered}
\end{equation}
The source operations preserve the displayed carrier. For addition,
if $r+s\ne0$ then at least one of $r,s$ is nonzero, so the joined
label is $1_L$. For multiplication, $rs\ne0$ forces both labels to
be $1_L$. The semiring laws and the two homomorphisms follow from
the ring laws and distributivity in $L$.

The units in this original semiring are exactly
\begin{equation}\tag{ZIL2}
 S^\times=\{(u,1_L):u\in R^\times\},\qquad
 e^2=e,\quad e\ne1_S,\quad e\notin S^\times.
\end{equation}
Indeed a product equal to $(1_R,1_L)$ has ring product $1_R$ and
label meet $1_L$, which forces both labels to equal $1_L$.
Conversely the ordinary ring inverse gives an inverse in $S$.
The inequalities in (ZIL2) use $0_R\ne1_R$.
In particular the note's integer $0$ is $e$, whereas its new symbol
$\tau$ is $0_S$. The Boolean subset $\{\tau,e\}$ has internal
multiplicative identity $e$; its inclusion into $S$ is not a
unit-preserving semiring map.

Here is the localization calculation without any cancellativity
assumption. Let $A$ be a commutative unital semiring and let
$c\in A$ satisfy $c^2=c$. The set $cA$ is a semiring with inherited
operations, zero $0_A$ and \emph{internal identity $c$}. Then
\begin{equation}\tag{ZIL3}
 A[c^{-1}]\simeq cA,\qquad a\longmapsto ca,
 \qquad [a,c^n]\longmapsto ca\quad(n\ge0).
\end{equation}
To prove this, multiplication by $c$ preserves addition and
multiplication because $c^2=c$, and sends $1_A$ to the internal
unit $c$. Its image of the denominator $c$ is that unit. Conversely,
if $f:A\to B$ is a zero- and unit-preserving homomorphism and
$f(c)$ is invertible, the equality $f(c)^2=f(c)$ gives $f(c)=1_B$.
The formula $g(ca)=f(a)$ is well defined: $ca=cb$ implies
$f(a)=f(ca)=f(cb)=f(b)$. It preserves both operations, zero and
the internal unit, and is forced by $f$. This proves the universal
property and the fraction formula. In particular the complete
equality relation of the localization is
$a\sim b\Longleftrightarrow ca=cb$.

Apply (ZIL3) to every labelled zero, without discarding any label:
\begin{equation}\tag{ZIL4}
\begin{gathered}
 S[z_\mu^{-1}]\simeq\downarrow\mu
     :=\{\nu\in L:\nu\le\mu\},\qquad
 1_{\downarrow\mu}=\mu,\quad0_{\downarrow\mu}=0_L,\\
 \ell_\mu:S\longrightarrow\downarrow\mu,
       \qquad\ell_\mu(r,\lambda)=\lambda\wedge\mu,\\
 (r,\lambda)\sim_\mu(s,\nu)
       \quad\Longleftrightarrow\quad
       \lambda\wedge\mu=\nu\wedge\mu.
\end{gathered}
\end{equation}
Indeed $z_\mu(r,\lambda)=z_{\mu\wedge\lambda}$, and every
$z_\nu$ with $\nu\le\mu$ occurs by choosing $(r,\lambda)=z_\nu$.
This identifies the entire corner semiring $z_\mu S$, its operations
and its internal identity. In particular
\begin{equation}\tag{ZIL5}
 S[e^{-1}]\simeq L,\quad [e,1]^{-1}=[e,1]=[1_S,1]=1_L,
 \qquad S[\tau^{-1}]\simeq\{0=1\}.
\end{equation}
Thus $e$ really is invertible in this localization, and its inverse
is exactly its own image, equal to the localized identity.
In the first localization $\tau$ is still $0_L\ne1_L$; it has
not become invertible. In the second localization the actual
absorbing zero is inverted and the entire semiring is trivial.
These are different operations.

The complete fibres, including amplitude loss, are
\begin{equation}\tag{ZIL6}
\begin{split}
 \ell_\mu^{-1}(\eta)
   ={}&\{z_\lambda:\lambda\wedge\mu=\eta\}\\
      &{}\cup
       \begin{cases}
        \{(r,1_L):r\in R\setminus\{0_R\}\},&\eta=\mu,\\
        \varnothing,&\eta\ne\mu,
       \end{cases}
       \qquad(\eta\le\mu),\\
 \sigma^{-1}(\lambda)
   ={}&\begin{cases}
       \{z_\lambda\},&\lambda\ne1_L,\\
       R\times\{1_L\},&\lambda=1_L.
      \end{cases}
\end{split}
\end{equation}
These follow by enumerating the two parts of the carrier in (ZIL1).
Thus supported-zero inversion loses every amplitude distinction in
the top fibre and preserves all original lattice labels. Inverting
a smaller label further identifies exactly the labels having the
same meet with it. Formula (ZIL6) also includes $\mu=0_L$.

For the full multiplicative set
$D_R=\{(r,1_L):r\in R\}$, including $e$ but excluding $\tau$,
one has
\begin{equation}\tag{ZIL7}
 D_R^{-1}S\simeq L,\qquad
 [(r,\lambda),(a,1_L)]\longmapsto\lambda.
\end{equation}
After $e$ is inverted, every member of $D_R$ already equals the
unit, by (ZIL4) with $\mu=1_L$. Conversely any map inverting
$D_R$ inverts $e$. The two universal properties give (ZIL7).
For $R=\mathbb Z$ and $L=\mathbb B$, this is exactly the note's
Theorem \texttt{thm:frac-collapse}; that theorem is correct.
Its symbol $\operatorname{Frac}(S)$ denotes this specific semiring
localization, not a field of fractions of the ring $\mathbb Z$.

In contrast, if $M\subset\mathbb Z\setminus\{0\}$ is a
multiplicative set containing $1$, then
\begin{equation}\tag{ZIL8}
 \{(m,1_L):m\in M\}^{-1}G_L(\mathbb Z)
       \simeq G_L(\mathbb Z_M),\qquad
 [(r,\lambda),(m,1_L)]\longmapsto(r/m,\lambda).
\end{equation}
Fraction equality on the left requires a $t\in M$ such that
$t m' r=t m r'$ and $\lambda=\lambda'$; multiplication by
top-labelled denominators never changes a label. These conditions
are exactly equality of the ring fractions and their labels.
This proves injectivity. Every top-labelled fraction and every
zero label in the right side is attained, proving surjectivity.
The formulas for fraction addition and multiplication agree with
(ZIL1), so the displayed bijection is an isomorphism. This proves
the full-lattice extension of the note's other localization theorem.

\subsection{The prime spectrum before and after this localization}

All primes of $S$ are the following two kinds:
\begin{equation}\tag{ZIL9}
\begin{gathered}
 P_J=\{z_\lambda:\lambda\in J\}
       \quad(J\in\operatorname{SpecLat}L),\\
 Q_{\mathfrak p}=\{z_\lambda:\lambda\in L\}
       \cup\{(r,1_L):r\in\mathfrak p\}
       \quad(\mathfrak p\in\operatorname{Spec}R),\\
 P_J\subsetneq Q_{\mathfrak p},\qquad
 D_S(z_\mu)=\{P_J:\mu\notin J\},\qquad
 V_S(e)=\{Q_{\mathfrak p}:\mathfrak p\in\operatorname{Spec}R\}.
\end{gathered}
\end{equation}
Here a lattice prime ideal is downward closed, closed under joins,
contains $0_L$, excludes $1_L$, and is prime for meets.
For completeness, a prime ideal of $S$ excluding $e$ contains no
top-labelled element: multiplication of any such element by $e$
is $e$. Its remaining labels form a proper lattice ideal, and
the primality test is exactly meet-primality. A prime containing
$e$ contains every $z_\lambda$, since $e z_\lambda=z_\lambda$.
Its top-labelled amplitudes form a ring ideal (additive negatives
are obtained by multiplying by $(-1_R,1_L)$); the prime condition
is ring primality. These arguments prove exhaustion and both
converses. Every $P_J$ is contained in every $Q_{\mathfrak p}$,
strictly because $e$ is in the latter and not the former.
Membership of $z_\mu$ now proves the last two formulas.

The spectrum of (ZIL4) identifies exactly with this supported open:
\begin{equation}\tag{ZIL10}
 \operatorname{Spec}(\downarrow\mu)
   \simeq\{J\in\operatorname{SpecLat}L:\mu\notin J\}
   \simeq D_S(z_\mu).
\end{equation}
Explicitly $J$ maps to $J\cap\downarrow\mu$. Conversely a prime
$H$ of $\downarrow\mu$ maps to
$\{\lambda:\lambda\wedge\mu\in H\}$.
These are lattice ideals and prime by meet and join preservation.
If $\mu\notin J$, primality gives
$\lambda\wedge\mu\in J\Longleftrightarrow\lambda\in J$,
so the maps are inverse. Basic opens correspond under the same
formula, proving the topological assertion.
Every arithmetic prime disappears under this localization because
every $Q_{\mathfrak p}$ contains $z_\mu$. This includes $\mu=1_L$;
it does not assert that these primes were absent from the original
geometry. For $L=\mathbb B$, $R=\mathbb Z$, the original strict
chain remains $\{\tau\}\subsetneq\{\tau,e\}\subsetneq Q_{(p)}$.

\subsection{The actual ASL coefficient map and its localization square}

Let $K$ be the tropical semifield with zero used in ASL18--22, and set
$K_L=K\otimes_{\mathbb B}L$. The zero and unit of $K$ are written
$\mathbf0_K,\mathbf1_K$. The nonzero indicator
$\beta_K:K\to\mathbb B$ is a semiring homomorphism: a maximum is
zero exactly when both inputs are zero, and a product is zero
exactly when a factor is zero. Consequently
\begin{equation}\tag{ZIL11}
\begin{gathered}
 \iota_L:L\longrightarrow K_L,\quad
       \iota_L(\lambda)=\mathbf1_K\otimes\lambda,\\
 \rho_L:K_L\longrightarrow L,\quad
 \rho_L\left(\sum_i a_i\otimes\lambda_i\right)
       =\bigvee_{i:a_i\ne\mathbf0_K}\lambda_i,\qquad
 \rho_L\iota_L=\operatorname{id}_L.
\end{gathered}
\end{equation}
The second map is $\beta_K\otimes\operatorname{id}_L$ followed
by the canonical isomorphism $\mathbb B\otimes_{\mathbb B}L=L$;
this proves independence from every tensor presentation. It also
proves the displayed formula and the retraction identity.
In particular $\iota_L$ is injective for arbitrary $L$, without
selecting a single Boolean character or replacing $L$ by one.

Retain the actual original-amplitude map
\begin{equation}\tag{ZIL12}
 \Phi_R:S\longrightarrow P_R:=R\times K_L,\qquad
       \Phi_R(r,\lambda)=(r,\iota_L(\lambda)).
\end{equation}
Both coordinates preserve the semiring operations, zero and unit.
If two images coincide their ring coordinates agree, and (ZIL11)
gives equality of their lattice coordinates. Thus $\Phi_R$ is
injective. Its image consists exactly of
$(r,\iota_L(\lambda))$ with $r\ne0_R$ implying $\lambda=1_L$.
The source's complete information is retained by this map.

For each $\mu\in L$, write
$d_\mu=\iota_L(\mu)$ and $u_\mu=(0_R,d_\mu)=\Phi_R(z_\mu)$.
The corresponding localization of the \emph{full product} is
\begin{equation}\tag{ZIL13}
 \begin{gathered}
 P_R[u_\mu^{-1}]\simeq d_\mu K_L,
       \qquad(r,w)\longmapsto d_\mu w,\\
 1_{d_\mu K_L}=d_\mu,\qquad
 P_R[\Phi_R(e)^{-1}]\simeq K_L,\quad(r,w)\longmapsto w.
 \end{gathered}
\end{equation}
Indeed $u_\mu^2=u_\mu$, and its corner in $P_R$ is
$\{(0_R,d_\mu w):w\in K_L\}$. Apply (ZIL3) and identify this
corner with its second coordinate. In particular the localization
of the product is not the zero ring when $\mu\ne0_L$; it is a
semiring whose entire ring factor has been erased.

Here is the exact square, including every label and both new units:
\begin{equation}\tag{ZIL14}
\begin{array}{ccc}
G_L(R)&\xrightarrow{\quad\Phi_R\quad}&R\times K_L\\[2pt]
{\scriptstyle(r,\lambda)\mapsto\lambda\wedge\mu}\big\downarrow
&&\big\downarrow{\scriptstyle(r,w)\mapsto d_\mu w}\\[2pt]
\downarrow\mu&\xrightarrow{\quad\eta\mapsto\iota_L(\eta)\quad}
   &d_\mu K_L .
\end{array}
\qquad
1_{\downarrow\mu}=\mu,quad1_{d_\mu K_L}=d_\mu.
\end{equation}
Both composites send $(r,\lambda)$ to
$\iota_L(\lambda\wedge\mu)$, because $\iota_L$ preserves meet.
The lower arrow takes values in the displayed corner since
$\eta\le\mu$ implies $d_\mu\iota_L(\eta)=\iota_L(\eta)$.
It preserves the new unit as well as zero and both operations.
Its retraction is the restriction of $\rho_L$, whose values lie in
$\downarrow\mu$ because
$\rho_L(d_\mu w)=\mu\wedge\rho_L(w)$.

The square is a pushout of unital semirings:
\begin{equation}\tag{ZIL15}
 (R\times K_L)\otimes_S S[z_\mu^{-1}]
       \simeq (R\times K_L)[u_\mu^{-1}]
       \simeq d_\mu K_L.
\end{equation}
To prove the asserted universal property directly, a map
$f:R\times K_L\to A$ agreeing on $S$ with a map from
$S[z_\mu^{-1}]$ sends $u_\mu$ to an invertible element.
By (ZIL3) it factors uniquely through $d_\mu K_L$.
Conversely any such factor gives the two agreeing maps in the
square. These constructions are inverse, proving (ZIL15).
This statement remains valid for $\mu=0_L$, when every receiving
semiring in this universal property is trivial.

At $\mu=1_L$, the complete relations erased by the two vertical
maps are respectively
\begin{equation}\tag{ZIL16}
 \begin{split}
 (r,\lambda)\sim_e(s,\nu)&\Longleftrightarrow\lambda=\nu,\\
 (r,w)\sim_{\Phi_R(e)}(s,v)&\Longleftrightarrow w=v.
 \end{split}
\end{equation}
The first relation is exactly the congruence generated by
$e\sim1_S$: in that congruence
$x=1_Sx\sim ex=z_{\sigma(x)}$, whereas $\sigma$ itself respects
$e\sim1_S$ and distinguishes every label. The second follows
in the same way from (ZIL3). Thus no amplitude dependence can
be recovered from either localized value alone.

Under the standing assumptions, the square at $\mu=1_L$ is not
a pullback. Its actual
pullback has the explicit form
\begin{equation}\tag{ZIL17}
 (R\times K_L)\times_{K_L}L\simeq R\times L,
 \quad((r,\iota_L(\lambda)),\lambda)\longleftrightarrow(r,\lambda).
\end{equation}
The original $S$ is precisely the subsemiring of this pullback
cut out by $r\ne0_R\Rightarrow\lambda=1_L$.
For example $(1_R,0_L)$ belongs to the pullback and not to $S$.
The formula for the pullback follows from equality of the second
coordinate with $\iota_L(\lambda)$, and its operations are
coordinatewise, proving the asserted isomorphism.

There is, however, a larger explicitly constructed subsemiring on
which the faithful map has a retraction:
\begin{equation}\tag{ZIL18}
 \begin{gathered}
 T_R=\{(r,w)\in R\times K_L:
             r\ne0_R\Longrightarrow\rho_L(w)=1_L\},\\
 \Psi_R:T_R\longrightarrow S,\quad
       \Psi_R(r,w)=(r,\rho_L(w)),\qquad
 \Psi_R\Phi_R=\operatorname{id}_S.
 \end{gathered}
\end{equation}
For closure under addition, a nonzero sum $r+s$ has at least
one nonzero summand, whose second-coordinate support is $1_L$;
the joined support is then $1_L$. A nonzero product $rs$ has
both factors nonzero, whose supports have meet $1_L$.
Zero and unit belong to $T_R$. Thus it is a subsemiring.
The map $\operatorname{id}_R\times\rho_L$ is a semiring
homomorphism on the entire product; its values lie in $S$
exactly on $T_R$, proving both the restriction and its maximal
domain for this formula. Formula (ZIL11) proves the retraction.
The localization of $T_R$ at $u_\mu$ is again $d_\mu K_L$,
because $T_R$ contains every $(0_R,w)$ and its $u_\mu$-corner
is exactly $\{0_R\}\times d_\mu K_L$.
After localization, $\Psi_R$ is the restricted $\rho_L$ in
(ZIL14). This gives exact compatible retractions before and
after localization on the stated domains.

No retraction $R\times K_L\to S$ of $\Phi_R$ exists on the
entire product. Suppose $F$ were one. Put $a=(1_R,\mathbf0_{K_L})$
and $b=(0_R,\mathbf1_{K_L})=\Phi_R(e)$. Then $ab=0$,
$a+b=1$, and $F(b)=e$. In $S$, the equality $ex=\tau$ forces
$x=\tau$ by (ZIL1). Hence $F(a)e=\tau$ forces $F(a)=\tau$,
and $1_S=F(a+b)=\tau+e=e$, contradicting (ZIL2).
This proves the exact obstruction and identifies the concrete
domain (ZIL18) on which the retraction does exist.

The ordinary ring-coordinate square has a different lower target:
\begin{equation}\tag{ZIL19}
\begin{array}{ccc}
 S&\xrightarrow{q}&R\\
 \big\downarrow&&\big\downarrow\\
 S[e^{-1}]\simeq L&\longrightarrow&R[0_R^{-1}]=\{0=1\}.
\end{array}
\end{equation}
Indeed $q(e)=0_R$, whose inverse forces $1_R=0_R$.
No map from $L$ to a nonzero ring can be a unital semiring
homomorphism: the equality $1_L+1_L=1_L$ would give $1+1=1$,
hence $1=0$, in the target ring. On the original $S$, however,
the combined map $(q,\sigma):S\hookrightarrow R\times L$
is injective. Equivalently the intersection of their equality
congruences is the diagonal. Thus keeping the unlocalized ring
coordinate alongside the full support coordinate recovers the
original element uniquely, whereas factoring that coordinate
through supported-zero inversion is impossible in a nonzero ring.

\subsection{Analytic zero locations and an exact non-detection test}

Let $\Omega$ be any set and $f:\Omega\to\mathbb C$ any function.
Form its \emph{top-labelled family} in $G_L(\mathbb C)$,
without asserting that the map $\mathbb C\to G_L(\mathbb C)$
preserves additive zero. Direct substitution gives
\begin{equation}\tag{ZIL20}
 \begin{gathered}
 \widehat f(z)=(f(z),1_L),\qquad
 \{z:f(z)=0\}=\widehat f^{-1}(\{e\}),\\
 \ell_\mu(\widehat f(z))=\mu\quad\hbox{for every }z,\qquad
 \Phi_{\mathbb C}(\widehat f(z))=(f(z),\mathbf1_{K_L}),\\
 d_\mu\,\operatorname{pr}_2\Phi_{\mathbb C}(\widehat f(z))
       =d_\mu\quad\hbox{for every }z.
 \end{gathered}
\end{equation}
Thus the right vertical map in (ZIL14) sends
$(f(z),\mathbf1_{K_L})$ to $d_\mu$ for every $z$;
at $\mu=1_L$ this is the constant unit $\mathbf1_{K_L}$.
The original family never takes the value $\tau$.
Its zero-amplitude locus is the inverse image of $e$, which
localization identifies with the unit. The ring projection of
the unlocalized joint family still recovers $f$ exactly.
These formulas apply to $\zeta$ on any domain avoiding its pole,
to $\xi$, or to a specified $H_t$ wherever they are defined.
No analytic theorem is needed for the calculation: it is valid
for every function and therefore does not distinguish their
different possible zero configurations.

The following finite calculation further determines exactly what
inversion alone can detect even when the backward heat equation
and a nonnegative real-root threshold are also retained. It is
a test of those stated pieces of data, not a replacement for the
zeta heat kernel. For $c\in\{0,1\}$ set
\begin{equation}\tag{ZIL21}
 P_t^{(c)}(z)=z^2+c-2t,\qquad
 \partial_tP_t^{(c)}=-2=-\partial_z^2P_t^{(c)},\qquad
 \{\hbox{all roots real}\}\Longleftrightarrow t\ge c/2.
\end{equation}
The differential identity is obtained by differentiating the
displayed polynomial. Its roots solve $z^2=2t-c$; for real $t$
they are both real, allowing the double root $0$, exactly when
$2t-c\ge0$. Thus the two thresholds are $0$ and $1/2$,
both nonnegative. At $t=0$ the first polynomial has all roots
real and the second has the nonreal roots $\pm i$.
Nevertheless every top-labelled family in (ZIL21) has the same
localized value in (ZIL20), for every $t,z,c$.
Consequently that localized value, the displayed heat equation,
and nonnegativity of the threshold do not together determine
whether the threshold is $0$ or strictly positive. This is a
counterexample to detection by those data only. It says nothing
contrary to any additional theorem about the actual $H_t$.

\subsection{Arithmetic ideal norms and the note's scaling error}

For $R=\mathbb Z$, retain the arithmetic ideals
$Q_{(n)}=\{z_\lambda:\lambda\in L\}
\cup\{(r,1_L):r\in n\mathbb Z\}$, for $n\ge1$.
The exact quotient homomorphism and its chosen norm are
\begin{equation}\tag{ZIL22}
 \pi_n:S\longrightarrow\mathbb Z/n\mathbb Z,\quad
       \pi_n(r,\lambda)=r\bmod n,\qquad
 \pi_n^{-1}(0)=Q_{(n)},\quad N(Q_{(n)})=n.
\end{equation}
The map is the ring projection followed by reduction modulo $n$,
so it is surjective and the displayed kernel is exact. The norm
here refers to this specific quotient congruence; an ideal alone
must not silently be substituted for a congruence.
For $\operatorname{Re}s>1$, this arithmetic family therefore gives
$\sum_{n\ge1}N(Q_{(n)})^{-s}=\sum_{n\ge1}n^{-s}$.
This proves the retained equality with the Dirichlet-series
definition of $\zeta(s)$ on that domain; no continuation or
zero-location conclusion follows from the equality alone.

Every ideal in (ZIL22) contains $e$ and hence $z_\mu$.
Its extension under $\ell_\mu$ is the entire $\downarrow\mu$,
because the image contains the new unit $\mu$.
Already $Q_{(2)}$ and $Q_{(3)}$ have different norms and identical
extended ideals. Thus there is no function of the extended ideal
alone recovering the original norm on this family. This is an
exact loss of arithmetic information in the localization, not
an assertion that the original prime or zeta data was unaffected.

The note contains a directly related false statement in its
class-semigroup definition. Write $A=\mathcal O_F$ for a number
field $F$, and $S_A=A\sqcup\{\tau\}=G_{\mathbb B}(A)$.
This notation avoids confusing the number field with the tropical
coefficient semifield $K$. With the source's own convention,
a nonzero semiring ideal means an ideal $I\ne\{\tau\}$.
Every such ideal is $I_J=J\cup\{\tau\}$ for an ideal
$J\subset A$, including $J=(0)$: multiplying any integer
member by the supported zero produces $e$, and the remaining
ideal rules give precisely the ordinary ring ideal. The note
defines
\begin{equation}\tag{ZIL23}
 I\sim J\quad\Longleftrightarrow\quad
 \exists a,b\in S_A\setminus\{\tau\}:aI=bJ.
\end{equation}
This allowed scalar set is $A$, which includes the supported
zero $e=0_A$. For every $I\ne\{\tau\}$,
\begin{equation}\tag{ZIL24}
 eI=\{\tau,e\},\qquad
 \{\hbox{nonzero ideals of }S_A\}/\!\sim
       \;=\;\{\hbox{one class}\}.
\end{equation}
Indeed $e\tau=\tau$, every ordinary member of $I$ is sent
to $e$, and a nonzero ideal has such a member. Taking $a=b=e$
in (ZIL23) then relates every pair. In particular the sentence
claiming these scalars are ``exactly the nonzero algebraic
integers'' is false. The note's subsequent injectivity proof
silently changes the scalar domain to $A\setminus\{0_A\}$,
and its surjectivity proof silently discards the allowed
nonzero semiring ideal $\{\tau,e\}=I_{(0)}$.

There is a precise corrected correspondence, with both restrictions
stated explicitly. Use only $I_J$ with $J\ne(0)$ and use only
scalars $a,b\in A\setminus\{0_A\}$. Then
\begin{equation}\tag{ZIL25}
 \begin{gathered}
 J\longleftrightarrow I_J=J\cup\{\tau\},\qquad
 I_{J_1}I_{J_2}=I_{J_1J_2},\\
 aI_{J_1}=bI_{J_2}\quad\Longleftrightarrow\quad
 aJ_1=bJ_2\qquad(a,b\ne0_A).
 \end{gathered}
\end{equation}
For the product identity, the ideal product consists of finite
sums of products. Terms involving $\tau$ contribute $\tau$;
the other products and sums give exactly the ordinary ideal
product, including its integer zero $e$. The scaling equivalence
is proved by intersecting with $A$ in one direction, and adjoining
$\tau$ in the other. These calculations give an exact isomorphism
with the usual nonzero integral-ideal class monoid defined by
nonzero scalar equivalence. Equivalence to fractional scaling is
also exact: $aJ_1=bJ_2$ means $J_2=(a/b)J_1$, and any element
of $F^\times$ has such a representation because $F$ is the
fraction field of $A$. Every nonzero fractional ideal has an integral
representative after multiplying by a common nonzero denominator.
Thus this restricted quotient is the usual ideal-class quotient
of fractional ideals by $F^\times$; in the Dedekind-domain setting
of the note it is the ordinary ideal class group.

If only the scalar restriction is repaired but $I_{(0)}$ is
retained, its class is an extra absorbing class. Nonzero scalars
cannot take a nonzero ideal of the domain $A$ to $(0)$, so this
class is distinct from all classes in (ZIL25); its product with
each of them is itself. Thus both source-domain corrections are
necessary for the claimed class-group statement, and the actual
definition as printed instead has the collapse (ZIL24).

\paragraph{Exact consequence for the proposed RH inference.}
The checked zero-inversion theorem says that the supported-zero
idempotent becomes the \emph{unit of a localized semiring}; it
does not give an inverse of $0_{\mathbb C}$ in the analytic
amplitude field. The complete map to the localized object is
(ZIL4), and its entire lost relation is (ZIL6). The coefficient
comparison retains the amplitudes before localization by (ZIL12)
and has the explicit connecting maps and retraction domain
(ZIL14)--(ZIL19). For an analytic family the exact zero locus
is retained on that unlocalized amplitude coordinate and is
erased by the localization, as (ZIL20) proves. The prime and
ideal-norm calculations (ZIL9)--(ZIL10), (ZIL22) identify precisely
which arithmetic information has been removed. These established
facts supply neither an off-critical zero nor a zero-location
contradiction for the actual zeta function. Formula (ZIL21)
refutes the narrower attempted inference from inversion, the
backward heat equation and threshold nonnegativity alone.
The calculation therefore constructs the exact resulting object
and its connection to the original one, without treating the
localization theorem or the note's erroneous scaling relation
as a disproof of RH.
